\chapter{FORMULÁRIO DE PESQUISA DE INFORMAÇÕES}
\label{ap:formulario}

Este apêndice apresenta o formulário completo utilizado para coletar dados sobre hábitos de uso de dispositivos móveis e percepções dos usuários. As informações coletadas foram fundamentais para o levantamento de requisitos e desenvolvimento das funcionalidades do aplicativo Desconecta.

\begin{questionframe}

\textbf{1. Você sente que passa mais tempo no celular do que gostaria?}\\
\phantom{\textbf{1. }}\opcao Sim, frequentemente\\
\phantom{\textbf{1. }}\opcao Às vezes\\
\phantom{\textbf{1. }}\opcao Não

\vspace{0.5cm}

\textbf{2. Com que frequência você pega o celular sem um motivo específico (por tédio, hábito, costume)?}\\
\phantom{\textbf{2. }}\opcao Quase o tempo todo\\
\phantom{\textbf{2. }}\opcao Às vezes\\
\phantom{\textbf{2. }}\opcao Raramente

\vspace{0.5cm}

\textbf{3. Você já sentiu ansiedade ou desconforto ao ficar longe do celular?}\\
\phantom{\textbf{3. }}\opcao Sim\\
\phantom{\textbf{3. }}\opcao Às vezes\\
\phantom{\textbf{3. }}\opcao Não

\vspace{0.5cm}

\textbf{4. Qual dessas situações acontece com você com mais frequência? (pode marcar mais de uma)} (múltipla escolha)\\
\phantom{\textbf{4. }}\opcao Ficar mexendo no celular até tarde da noite\\
\phantom{\textbf{4. }}\opcao Olhar o celular logo ao acordar\\
\phantom{\textbf{4. }}\opcao Se distrair com notificações no trabalho ou nos estudos\\
\phantom{\textbf{4. }}\opcao Usar o celular mesmo com pessoas por perto\\
\phantom{\textbf{4. }}\opcao Nenhuma dessas

\vspace{0.5cm}

\textbf{5. Você já tentou reduzir o tempo de uso do celular por conta própria?}\\
\phantom{\textbf{5. }}\opcao Sim, mas não consegui manter\\
\phantom{\textbf{5. }}\opcao Sim, e consegui\\
\phantom{\textbf{5. }}\opcao Não, nunca tentei

\vspace{0.5cm}

\textbf{6. O que mais te motivaria a usar menos o celular?} (múltipla escolha)\\
\phantom{\textbf{6. }}\opcao Ter mais tempo livre\\
\phantom{\textbf{6. }}\opcao Dormir melhor\\
\phantom{\textbf{6. }}\opcao Melhorar o foco e a produtividade\\
\phantom{\textbf{6. }}\opcao Reduzir ansiedade e estresse\\
\phantom{\textbf{6. }}\opcao Outro: \rule{7cm}{0.4pt}

\vspace{0.7cm}

\textbf{7. Você acha que participar de um desafio em grupo focado na redução do tempo de tela pode ajudar você a usar menos o celular?}\\
\phantom{\textbf{7. }}\opcao Sim\\
\phantom{\textbf{7. }}\opcao Talvez\\
\phantom{\textbf{7. }}\opcao Não

\vspace{0.5cm}

\textbf{8. Ver outras pessoas compartilhando o próprio tempo de tela te motivaria a reduzir o seu também?}\\
\phantom{\textbf{8. }}\opcao Sim, bastante\\
\phantom{\textbf{8. }}\opcao Um pouco\\
\phantom{\textbf{8. }}\opcao Não faria diferença

\vspace{0.5cm}

\textbf{9. Você se sente mais motivado(a) quando está em um grupo com o mesmo objetivo?}\\
\phantom{\textbf{9. }}\opcao Sim\\
\phantom{\textbf{9. }}\opcao Às vezes\\
\phantom{\textbf{9. }}\opcao Não

\vspace{0.5cm}

\textbf{10. Saber que outras pessoas podem ver o seu tempo de tela faria você usar menos o celular?}\\
\phantom{\textbf{10. }}\opcao Sim\\
\phantom{\textbf{10. }}\opcao Talvez\\
\phantom{\textbf{10. }}\opcao Não\\
\phantom{\textbf{10. }}\opcao Não sei

\vspace{0.5cm}

\textbf{11. Você concorda com a ideia de participar de um grupo onde as pessoas se apoiam e compartilham a redução diária do tempo de tela?}\\
\phantom{\textbf{11. }}\opcao Sim\\
\phantom{\textbf{11. }}\opcao Talvez\\
\phantom{\textbf{11. }}\opcao Não

\vspace{0.5cm}

\textbf{12. O que mais te ajudaria em um grupo focado em usar menos o celular?} (múltipla escolha)\\
\phantom{\textbf{12. }}\opcao Ver o progresso de outras pessoas\\
\phantom{\textbf{12. }}\opcao Receber incentivo e lembretes\\
\phantom{\textbf{12. }}\opcao Ter pequenas metas diárias\\
\phantom{\textbf{12. }}\opcao Competição leve (ranking, pontos, medalhas)\\
\phantom{\textbf{12. }}\opcao Só acompanhar, sem participar ativamente\\
\phantom{\textbf{12. }}\opcao Outro: \rule{7cm}{0.4pt}

\end{questionframe}

\vspace{1cm}

\textbf{Observações:} Este formulário foi aplicado de forma online e as respostas foram coletadas de forma anônima, garantindo a privacidade dos participantes. Os dados coletados foram utilizados exclusivamente para fins acadêmicos no desenvolvimento deste trabalho.